Vi har taget udgangspunkt ud fra en artikel der specifikt undersøger hvordan datacentre påvirker miljøet og hvilket strategier man kan anvende, vi vurderer denne kilde værdig da den er skrevet af en gruppe af professorer og eksperter inden for Micro Electronics/Research og Environmental \& Reliability Engineering 

Artikklen er skrevet i 2025 hvilket gør den mere relevant 

Brugen af energi er stigende og ud fra ekspertviden ligger der allerede en forventning på 20\% stigning i energiforbrug på datacentre samt en 13\% stigning i drivhusgas udledning bare på datacentre alene estimeret til 2030. Ca. 1,4\% af vores drivhusgas i 2020 blev udlet grundet kommunikation og informationsteknologi, og imellem 2010 – 2018 var der en 6\% stigning i energiforbrug på datacentre. Pga. effektivitets og forbrugs stigninger på hardwaren og software. 

Største miljømæssige problemer inden for app branchen kommer fra energien som datacentrene anvender, nogle datacentre anvender vandkraft, og strøm fremproduceret af fossile brandstoffer såsom kul. 

I vores overvejelse har vi herefter lavet overvejelser hvordan vores produkt også kan understøtte klimaet mest mulig og har derfor lavet en tabel over hvilket energikilder vi vurderer som mest miljøvenlige samt realistisk at indføre i datacentre 


\begin{table}[H]
    \centering
    \begin{tabular}{|l|l|}
        \hline
        \textbf{Energikilde} & \textbf{Rangering} \\ \hline
        Vindenergi          & 1 \\ \hline
        Solenergi           & 2 \\ \hline
        Vandkraft           & 3 \\ \hline
        Biomasse            & 4 \\ \hline
        Atomkraft           & 5 \\ \hline
    \end{tabular}
    \caption{Rangering af energikilder efter miljøpåvirkning\footfullcite{LifeCycleGreenhouseGas2021}}
    \label{tab:co2_udledning}
\end{table}


Vi har valgt Vindenergi solenergi og vandkraft i toppen da det er en af de mindst c02 udledende energikilder som der kan findes, samt hvis man bruger nedkøling som er baseret på koldt miljø såsom at anvende Grønlands naturligt kølige klima til at bygge datacentre som anvender naturens klima, til at køle serverne ned frem for at brug en ekstern energikilde til at drive kølerne således at vi minimerer forbrug per datacentre. 

Problemet er at denne løsning ikke er 100\% fremtidssikret samt vedligeholdelse af centrene bliver besværlige, grundet befolkningstallet i Grønland samt at få en person derned til vedligeholdelse. 